\documentclass[10pt,a4paper]{amsart}
\usepackage{amsmath,amssymb,amsthm,graphicx}
% Retained minimal preamble for the Stacks Project chapter more-algebra.tex (More on Algebra).
% The chapter file pinned-original-source.tex carries only \input{preamble} on its first line, so no
% preamble text lives inside the chapter itself.  The Stacks Project uses one single shared preamble for
% every chapter; the snapshot stages it identically as sources/stacks-chow/preamble.tex and
% sources/stacks-etale-cohomology/preamble.tex.  The block below is copied mechanically, line for line,
% from sources/stacks-chow/preamble.tex lines 14-18 (the Xy-pic setup) and 154-217 (the native theorem
% declarations together with the mathematical macros).  Nothing was added, removed or redefined; the
% wrapper class and the amsmath/amssymb/amsthm/graphicx base packages are supplied by the assembler.
% Chapter: more-algebra.tex -- More on Algebra.
% For commutative diagrams we use Xy-pic
\usepackage[all]{xy}

% We use 2cell for 2-commutative diagrams.
\xyoption{2cell}

% Theorem environments.
%
\theoremstyle{plain}
\newtheorem{theorem}[subsection]{Theorem}
\newtheorem{proposition}[subsection]{Proposition}
\newtheorem{lemma}[subsection]{Lemma}

\theoremstyle{definition}
\newtheorem{definition}[subsection]{Definition}
\newtheorem{example}[subsection]{Example}
\newtheorem{exercise}[subsection]{Exercise}
\newtheorem{situation}[subsection]{Situation}

\theoremstyle{remark}
\newtheorem{remark}[subsection]{Remark}
\newtheorem{remarks}[subsection]{Remarks}

\numberwithin{equation}{subsection}

% Macros
%
\def\lim{\mathop{\mathrm{lim}}\nolimits}
\def\colim{\mathop{\mathrm{colim}}\nolimits}
\def\Spec{\mathop{\mathrm{Spec}}}
\def\Hom{\mathop{\mathrm{Hom}}\nolimits}
\def\Ext{\mathop{\mathrm{Ext}}\nolimits}
\def\SheafHom{\mathop{\mathcal{H}\!\mathit{om}}\nolimits}
\def\SheafExt{\mathop{\mathcal{E}\!\mathit{xt}}\nolimits}
\def\Sch{\mathit{Sch}}
\def\Mor{\mathop{\mathrm{Mor}}\nolimits}
\def\Ob{\mathop{\mathrm{Ob}}\nolimits}
\def\Sh{\mathop{\mathit{Sh}}\nolimits}
\def\NL{\mathop{N\!L}\nolimits}
\def\CH{\mathop{\mathrm{CH}}\nolimits}
\def\proetale{{pro\text{-}\acute{e}tale}}
\def\etale{{\acute{e}tale}}
\def\QCoh{\mathit{QCoh}}
\def\Ker{\mathop{\mathrm{Ker}}}
\def\Im{\mathop{\mathrm{Im}}}
\def\Coker{\mathop{\mathrm{Coker}}}
\def\Coim{\mathop{\mathrm{Coim}}}

% Boxtimes
%
\DeclareMathSymbol{\boxtimes}{\mathbin}{AMSa}{"02}

%
% Macros for moduli stacks/spaces
%
\def\QCohstack{\mathcal{QC}\!\mathit{oh}}
\def\Cohstack{\mathcal{C}\!\mathit{oh}}
\def\Spacesstack{\mathcal{S}\!\mathit{paces}}
\def\Quotfunctor{\mathrm{Quot}}
\def\Hilbfunctor{\mathrm{Hilb}}
\def\Curvesstack{\mathcal{C}\!\mathit{urves}}
\def\Polarizedstack{\mathcal{P}\!\mathit{olarized}}
\def\Complexesstack{\mathcal{C}\!\mathit{omplexes}}
% \Pic is the operator that assigns to X its picard group, usage \Pic(X)
% \Picardstack_{X/B} denotes the Picard stack of X over B
% \Picardfunctor_{X/B} denotes the Picard functor of X over B
\def\Pic{\mathop{\mathrm{Pic}}\nolimits}
\def\Picardstack{\mathcal{P}\!\mathit{ic}}
\def\Picardfunctor{\mathrm{Pic}}
\def\Deformationcategory{\mathcal{D}\!\mathit{ef}}
\usepackage{cleveref}
\crefname{originallemma}{original lemma}{original lemmas}
\Crefname{originallemma}{Original lemma}{Original lemmas}
\makeatletter
\newlabel{more-algebra-lemma-P-universally-injective}{{05GP}{source}{}{sourcepointer-1}{}}
\newlabel{more-algebra-lemma-P-universally-injective@cref}{{[originallemma][0][]05GP}{[1][0][]source}{}{}{}}
\makeatother
\crefname{originallemma}{original lemma}{original lemmas}
\Crefname{originallemma}{Original lemma}{Original lemmas}
\makeatletter
\newlabel{lemma-auto-ass-implies-P}{{05GN}{source}{}{sourcepointer-2}{}}
\newlabel{lemma-auto-ass-implies-P@cref}{{[originallemma][0][]05GN}{[1][0][]source}{}{}{}}
\makeatother
\title{A universally injective ring map from a P-ring condition}
\author{The Stacks Project Authors}
\date{stacks-project:more-algebra.tex:local-snapshot-20261012; GNU FDL 1.2 or later}
\begin{document}
\maketitle
\setcounter{subsection}{1}
\noindent\textit{Editorial note.} All delivered mathematical text is the author's own, copied from the pinned Stacks Project chapter file more-algebra.tex; no statement or proof step was rewritten, reordered or removed.\par\medskip
\noindent\textit{Editorial note.} The chapter file pinned-original-source.tex carries only \textbackslash{}input\{preamble\} on its first line, so no preamble text lives inside the chapter. The retained file sources/stacks-more-algebra/retained-preamble.tex is a mechanical, line-for-line copy of the project's single shared preamble as staged in the snapshot (Xy-pic setup plus theorem declarations and mathematical macros). Nothing was added, removed or redefined, and the excerpt serves only as preamble, never as a span coordinate.\par\medskip
\noindent\textit{Editorial note.} Both counted spans start at the \textbackslash{}begin line of their environment, so the statement and the author's proof are delivered whole; surrounding prose of the section is deliberately outside the counted spans and is carried as standard prerequisites instead.\par\medskip
\noindent\textit{Editorial note.} The citation fields (source\_id, source\_version, authors, source\_url, license.name, source\_title) are identical to those of every other card this lane has produced from this chapter, so the batch selector treats them as one source.\par\medskip
\noindent\textit{Allowed standard prerequisites (editorial).}\begin{itemize}
\item Absolutely integrally closed ring and the notion of a universally injective ring map. defined in the same section, source lines 3650-3707, before the counted statement. Supplies both the hypothesis and the property being proved.
\item A base change of a ring map can be written as a filtered colimit of finitely generated subalgebras. used at proof lines 3720-3740 to reduce the verification to the finite type case. Reduces an arbitrary base change to a manageable one.
\item Injectivity is preserved under filtered colimits of maps. used at proof lines 3740-3754. Transports the finite type verification back to the arbitrary base change.
\end{itemize}\medskip
\ifdefined\hypertarget\hypertarget{sourcepointer-1}{}\fi
\noindent\textit{Source reference.} Original lemma 05GP: focused-seeds/AGP/sources/stacks-snapshot-20261011/tags.txt, source lines 7068--7068. The lemma's own canonical tag, quoted so that readers can resolve the label inside the delivered text; it points at the counted statement itself and carries no proof content.\par\medskip
\ifdefined\hypertarget\hypertarget{sourcepointer-2}{}\fi
\noindent\textit{Source reference.} Original lemma 05GN: focused-seeds/AGP/sources/stacks-snapshot-20261011/tags.txt, source lines 7067--7067. Navigation pointer to a statement the counted proof refers to but does not contain; nothing load-bearing is delegated to it.\par\medskip
\begin{lemma}
\label{lemma-P-universally-injective}
Let $R$ be a ring having property (P) of
Lemma \ref{lemma-auto-ass-implies-P}.
Let $u : N \to M$ be a homomorphism of projective $R$-modules.
Then $u$ is universally injective if and only if $u$ is injective.
\end{lemma}

\begin{proof}
Assume $u$ is injective. Our goal is to show $u$ is universally injective.
First we choose a module $Q$ such that $N \oplus Q$ is free. On considering
the map $N \oplus Q \to M \oplus Q$ we see that it suffices to prove
the lemma in case $N$ is free. In this case $N$ is a directed colimit of
finite free $R$-modules. Thus we reduce to the case that $N$ is a finite
free $R$-module, say $N = R^{\oplus n}$. We prove the lemma by induction
on $n$. The case $n = 0$ is trivial.

\medskip\noindent
Let $u : R^{\oplus n} \to M$ be an injective module map with $M$ projective.
Choose an $R$-module $Q$ such that $M \oplus Q$ is free. After replacing
$u$ by the composition $R^{\oplus n} \to M \to M \oplus Q$ we see that
we may assume that $M$ is free. Then we can find a direct summand
$R^{\oplus m} \subset M$ such that $u(R^{\oplus n}) \subset R^{\oplus m}$.
Hence we may assume that $M = R^{\oplus m}$.
In this case $u$ is given by a matrix $A = (a_{ij})$ so that
$u(x_1, \ldots, x_n) = (\sum x_i a_{i1}, \ldots, \sum x_i a_{im})$.
As $u$ is injective, in particular
$u(x, 0, \ldots, 0) = (xa_{11}, xa_{12}, \ldots, xa_{1m}) \not = 0$ if
$x \not = 0$, and as $R$ has property (P) we see that
$a_{11}R + a_{12}R + \ldots + a_{1m}R = R$. Hence see that
$R(a_{11}, \ldots, a_{1m}) \subset R^{\oplus m}$ is a direct summand
of $R^{\oplus m}$, in particular $R^{\oplus m}/R(a_{11}, \ldots, a_{1m})$
is a projective $R$-module. We get a commutative diagram
$$
\xymatrix{
0 \ar[r] &
R \ar[rr] \ar[d]^1 & & R^{\oplus n} \ar[r] \ar[d]^u &
R^{\oplus n - 1} \ar[r] \ar[d] & 0 \\
0 \ar[r] & R \ar[rr]^{(a_{11}, \ldots, a_{1m})} & &
R^{\oplus m} \ar[r] & R^{\oplus m}/R(a_{11}, \ldots, a_{1m}) \ar[r] & 0
}
$$
with split exact rows. Thus the right vertical arrow is injective
and we may apply the induction hypothesis to conclude that
the right vertical arrow is universally injective. It follows that the
middle vertical arrow is universally injective.
\end{proof}

\end{document}
